\thispagestyle{empty}
\section*{Post Scriptum: Nada nuevo bajo el Sol}
\addcontentsline{toc}{section}{Post Scriptum}

Durante la gestación de este trabajo, la estructura matemática aquí presentada fue construida desde cero, guiada exclusivamente por la necesidad de capturar las dinámicas del Go. Casi al final noté algo inesperado, en lo esencial, coincide con lo que en la literatura se conoce como \textit{Teoría de Grafos Etiquetados}.

En esa teoría es habitual considerar un grafo $(G,A)$ junto con una función de etiquetado $s:G\to E$. En esta monografía, en cambio, representé los estados mediante el conjunto de pares $(E\times G)$, dando lugar a la estructura $(G,E,A)$ o, equivalentemente, al grafo etiquetado $(G,A,E)$.

En un primer momento consideré reescribir toda la obra para adoptar la notación estándar. Sin embargo, comprendí que ello no aportaría ninguna ventaja conceptual relevante. Desde la teoría de conjuntos, una función no es más que un conjunto particular de pares ordenados; ambas formulaciones contienen esencialmente la misma información.

Más importante aún, durante todo el desarrollo nunca fue necesario invocar propiedades específicas de la teoría de grafos. La construcción se realizó exclusivamente a partir de conjuntos, relaciones y funciones surgidas de la propia dinámica del juego. Mantener esa perspectiva permitió avanzar con mayor naturalidad, sin introducir herramientas externas que nunca resultaron relevantes.

Lejos de restar originalidad, esta coincidencia resulta profundamente satisfactoria. Partiendo únicamente de las reglas del Go, sin buscar deliberadamente una teoría preexistente, el proceso de abstracción condujo de manera natural a una estructura ya conocida por la matemática desde otra perspectiva.

\noindent
Buscar la esencia del Go fue, al final, encontrar una estructura que siempre había estado allí.
\vfill
\begin{flushright}
	\textit{Arias, W. Adrián}
\end{flushright}